\section{Positive \texorpdfstring{$\Ein$}{Ein} moment systems and
algebraically independent Jacobi spectra}
\label{sec:positive-ein-jacobi}

The generic Hankel system built directly from the moments
$\Ein(m+1)$ is quasi-definite but not positive-definite.  A finite
difference removes this defect and produces a genuine Hausdorff moment
problem without sacrificing any algebraic independence.

Fix $a,b\in\overline{\mathbb Q}_{>0}$ and an integer $r\geq1$.  For a
sequence $f=(f_m)_{m\geq0}$ write
\[
 (\mathsf Df)_m=f_{m+1}-f_m,
\]
and put
\begin{equation}\label{eq:positive-ein-moments}
 f_m=\Ein(b+am),\qquad
 \mu_m^{(a,b,r)}=(-1)^{r-1}(\mathsf D^rf)_m\quad(m\geq0).
\end{equation}
When there is no danger of confusion, write simply $\mu_m$.

\begin{theorem}[Positive finite-difference $\Ein$ moment theorem]
\label{thm:positive-ein-jacobi}
Let $K=\Eexp$.  The moments in
\eqref{eq:positive-ein-moments} have the exact representation
\begin{align}
 \mu_m
 &=\int_0^1 e^{-(b+am)t}\frac{(1-e^{-at})^r}{t}\,dt
 \label{eq:positive-ein-laplace}\\
 &=\int_{e^{-a}}^1 x^m w_{a,b,r}(x)\,dx,
 \qquad
 w_{a,b,r}(x)=
 \frac{x^{b/a-1}(1-x)^r}{-\log x}.
 \label{eq:positive-ein-weight}
\end{align}
In particular $(\mu_m)_{m\geq0}$ is a strictly positive-definite
Hausdorff moment sequence.  Moreover:
\begin{enumerate}
\item The countable family $\mu_0,\mu_1,\ldots$ is algebraically
independent over $K$, in the finite-subset sense.  The infinite Hankel
matrix $(\mu_{i+j})_{i,j\geq0}$ is strictly totally positive: every one
of its finite minors is positive.
\item With
\begin{equation}\label{eq:positive-ein-Hankel}
 \mathcal H_n=\det(\mu_{i+j})_{0\leq i,j\leq n},
 \qquad \mathcal H_{-1}=1,
\end{equation}
one has
\begin{equation}\label{eq:positive-ein-Heine}
 \mathcal H_n=
 \frac1{(n+1)!}
 \int_{[e^{-a},1]^{n+1}}
 \prod_{0\leq i<j\leq n}(x_i-x_j)^2
 \prod_{j=0}^n w_{a,b,r}(x_j)\,dx_j>0.
\end{equation}
The family $\mathcal H_0,\mathcal H_1,\ldots$ is algebraically
independent over $K$, and its exact capacity scale is
\begin{equation}\label{eq:positive-ein-Hankel-capacity}
 \lim_{n\to\infty}\mathcal H_n^{1/n^2}
 =\operatorname{cap}([e^{-a},1])=\frac{1-e^{-a}}4.
\end{equation}
\item Let $p_{-1}=0$, $p_0=1$ be the monic orthogonal polynomials for
the functional $\mathcal L(x^m)=\mu_m$, and normalize the Favard
recurrence by
\begin{equation}\label{eq:positive-ein-Favard}
 p_{n+1}(x)=(x-\alpha_n)p_n(x)-\beta_np_{n-1}(x).
\end{equation}
Then $\beta_n>0$ for every $n\geq1$, and
\[
 \mu_0,\alpha_0,\alpha_1,\ldots,
 \beta_1,\beta_2,\ldots
\]
is algebraically independent over $K$.
\item The Jacobi matrix
\begin{equation}\label{eq:positive-ein-Jacobi}
 J=
 \begin{pmatrix}
 \alpha_0&\sqrt{\beta_1}&&\\
 \sqrt{\beta_1}&\alpha_1&\sqrt{\beta_2}&\\
 &\sqrt{\beta_2}&\alpha_2&\ddots\\
 &&\ddots&\ddots
 \end{pmatrix}
\end{equation}
defines a bounded self-adjoint operator on $\ell^2(\mathbb N_0)$.
It has simple purely absolutely continuous spectrum
\[
 \sigma(J)=[e^{-a},1]
\]
of multiplicity one.
Its monic recurrence coefficients lie in the Nevai class of this interval:
\begin{equation}\label{eq:positive-ein-Nevai}
 \alpha_n\longrightarrow\frac{1+e^{-a}}2,
 \qquad
 \beta_n\longrightarrow\frac{(1-e^{-a})^2}{16}.
\end{equation}
\item For every $N\geq1$, the $N$ eigenvalues of the leading
$N\times N$ principal truncation $J_N$ are distinct, belong to
$(e^{-a},1)$, and are jointly algebraically independent over $K$.
If they are denoted in increasing order by
$\lambda_{N,1},\ldots,\lambda_{N,N}$, then their normalized counting
measures satisfy
\begin{equation}\label{eq:positive-ein-arcsine}
 \frac1N\sum_{j=1}^N\delta_{\lambda_{N,j}}
 \xrightarrow[N\to\infty]{\mathrm{weak}}
 \frac{\mathbf 1_{(e^{-a},1)}(x)}
 {\pi\sqrt{(x-e^{-a})(1-x)}}\,dx.
\end{equation}
\end{enumerate}
\end{theorem}

\begin{proof}
The standard integral
\[
 \Ein(s)=\int_0^1\frac{1-e^{-st}}t\,dt\qquad(s>0)
\]
shows, after taking the $r$th forward difference in $m$, that
\[
 (-1)^{r-1}(\mathsf D^rf)_m
 =\int_0^1e^{-(b+am)t}\frac{(1-e^{-at})^r}{t}\,dt.
\]
The substitution $x=e^{-at}$ gives
\eqref{eq:positive-ein-weight}.  The density is positive on
$(e^{-a},1)$.  At $x=1$ it is asymptotic to
$x^{b/a-1}(1-x)^{r-1}$, while the other endpoint is strictly positive;
hence the density is integrable and has essential support
$[e^{-a},1]$.  The Gram determinant formula now gives
\eqref{eq:positive-ein-Heine}; its integrand is positive off the
diagonals, so every $\mathcal H_n$ is strictly positive.

More generally, choose strictly increasing nonnegative integer lists
$i_1<\cdots<i_k$ and $j_1<\cdots<j_k$.  Andreief's identity writes the
corresponding Hankel minor as
\[
 \frac1{k!}\int_{[e^{-a},1]^k}
 \det(x_\ell^{i_p})_{p,\ell=1}^k
 \det(x_\ell^{j_p})_{p,\ell=1}^k
 \prod_{\ell=1}^kw_{a,b,r}(x_\ell)\,dx_\ell.
\]
On the ordered simplex $e^{-a}<x_1<\cdots<x_k<1$, each generalized
Vandermonde determinant is a positive Vandermonde times a Schur
polynomial with positive values.  Their product is positive on a set of
positive measure, proving strict total positivity.

For algebraic independence, let
\[
 y_m=\Ein(b+am).
\]
The points $b+am$ are distinct nonzero algebraic numbers, so the
multipoint $\Ein$ theorem makes every finite subfamily of the $y_m$
algebraically independent over $K$.  For fixed $M$, the $M+1$ linear
forms $\mu_0,\ldots,\mu_M$ in
$y_0,\ldots,y_{M+r}$ are linearly independent: the coefficient of
$y_{M+r}$ occurs only in $\mu_M$, and descending induction finishes the
argument.  Extending these coefficient vectors to a basis gives an
invertible linear change of variables.  Thus
$\mu_0,\ldots,\mu_M$ are algebraically independent over $K$.

Expansion of the determinant at its lower-right entry gives
\begin{equation}\label{eq:positive-ein-Hankel-triangular}
 \mathcal H_n=\mathcal H_{n-1}\mu_{2n}
 +R_n(\mu_0,\ldots,\mu_{2n-1}).
\end{equation}
Induction in $n$, viewing a putative relation as a polynomial in
$\mu_{2n}$, proves the algebraic independence of the $\mathcal H_n$.

Strict positivity of all Gram determinants gives the monic Favard
recurrence, with
\begin{equation}\label{eq:positive-ein-beta-Hankel}
 \beta_n=\frac{\mathcal H_n\mathcal H_{n-2}}
 {\mathcal H_{n-1}^2}>0.
\end{equation}
For each $N\geq0$, finite Gram--Schmidt elimination gives a rational map
\[
 (\mu_0,\ldots,\mu_{2N+1})
 \longmapsto
 (\mu_0,\alpha_0,\ldots,\alpha_N,
             \beta_1,\ldots,\beta_N).
\]
Conversely, the weighted Motzkin-path expansion, equivalently
$\mu_m=\mu_0(J^m)_{00}$, recovers the moments through order $2N+1$
polynomially from the displayed Favard variables.  The two lists have
$2N+2$ entries and the transformations are inverse on the open set of
nonzero Hankel determinants.  They are therefore birational.  Algebraic
independence of the moments proves that of the Favard variables.

Normalize the measure in \eqref{eq:positive-ein-weight} by its mass
$\mu_0$.  Polynomials are dense in its $L^2$ space, and the orthonormal
polynomial transform identifies $J$ with multiplication by $x$ on
\[
 L^2\!\left([e^{-a},1],\mu_0^{-1}w_{a,b,r}(x)\,dx\right).
\]
This proves bounded self-adjointness.  Since the density is positive
almost everywhere on the whole interval, multiplication by $x$ has
purely absolutely continuous spectrum $[e^{-a},1]$.  The constant
function is cyclic, so the spectral multiplicity is one.  The
Denisov--Rakhmanov theorem~\cite{Denisov2004}, affinely rescaled from $[-2,2]$ to
$[e^{-a},1]$, gives
\[
 \alpha_n\to\frac{1+e^{-a}}2,
 \qquad
 \sqrt{\beta_n}\to\frac{1-e^{-a}}4,
\]
which is exactly \eqref{eq:positive-ein-Nevai} in the monic
normalization \eqref{eq:positive-ein-Favard}.
If $h_n=\mathcal L(p_n^2)$, then
$h_n/h_{n-1}=\beta_n$ and
$\mathcal H_n=\prod_{j=0}^nh_j$.  Consequently
\[
 \mathcal H_n
 =\mu_0^{n+1}\prod_{k=1}^n\beta_k^{n+1-k}.
\]
Taking logarithms and using \eqref{eq:positive-ein-Nevai} proves
\eqref{eq:positive-ein-Hankel-capacity}.

It remains to prove the arithmetic assertion about finite truncations.
The positivity of the off-diagonal entries makes every $J_N$ an
irreducible real symmetric tridiagonal matrix.  Hence its eigenvalues are
simple; the standard zero theorem for orthogonal polynomials places them
in $(e^{-a},1)$.  Write the characteristic polynomial as
\[
 \det(TI_N-J_N)=T^N+c_1T^{N-1}+\cdots+c_N.
\]
Each $c_j$ is a polynomial over $\mathbb Q$ in
$\alpha_0,\ldots,\alpha_{N-1},\beta_1,\ldots,\beta_{N-1}$.
These $N$ polynomials are algebraically independent: after formally
specializing all $\beta_i$ to zero, they become the signed elementary
symmetric polynomials in the independent diagonal variables
$\alpha_0,\ldots,\alpha_{N-1}$.  The algebraic independence of the
actual Favard variables makes their evaluation map injective, so the
actual $c_1,\ldots,c_N$ are algebraically independent over $K$.

If $\lambda_{N,1},\ldots,\lambda_{N,N}$ are the eigenvalues, then the
$c_j$ are their elementary symmetric functions and the splitting field
is finite over $K(c_1,\ldots,c_N)$.  Consequently
\[
 \operatorname{trdeg}_K
 K(\lambda_{N,1},\ldots,\lambda_{N,N})=N.
\]
There are exactly $N$ eigenvalues, so they are jointly algebraically
independent over $K$.

Finally, these eigenvalues are the zeros of $p_N$.  Formula
\eqref{eq:positive-ein-Nevai} gives their limiting distribution directly.
Indeed, put
\[
 c=\frac{1+e^{-a}}2,
 \qquad d=\frac{1-e^{-a}}4.
\]
For every fixed integer $k\geq0$, a closed-walk expansion of the trace,
with the $O(k)$ boundary rows discarded, gives
\[
 \lim_{N\to\infty}\frac1N\operatorname{tr}(J_N^k)
 =\frac1{2\pi}\int_0^{2\pi}(c+2d\cos\theta)^k\,d\theta.
\]
The matrices $J_N$ are uniformly bounded, so polynomial approximation
extends this convergence from moments to every continuous test function.
The push-forward of $d\theta/(2\pi)$ under
$\theta\mapsto c+2d\cos\theta$ is precisely the probability measure on
the right side of \eqref{eq:positive-ein-arcsine}.  This proves the weak
convergence.  Equivalently, it is the classical zero-counting theorem for
a regular measure on a single interval.
\end{proof}

\begin{remark}[What is new and what is imported]
The Hausdorff moment representation and the spectral consequences are
classical once \eqref{eq:positive-ein-weight} is known.  The arithmetic
content is that the explicit moments, all Hankel determinants, all
Favard coordinates, and every finite Jacobi spectrum acquire the stated
algebraic independence from the multipoint $\Ein$ theorem.  Historical
priority for this exact arithmetic realization should be audited before
publication.
\end{remark}
